%------------------------------------------------------------------------------%
\chapter{Referências e Recursos Online}
\label{cap:Recursos_Online}

\minitoc

%------------------------------------------------------------------------------%
\section{Recursos Online}

Existem muitos recursos \textit{online} que servem como apoio ao estudo de Matemática e Cálculo.
Utilizar esses recursos é altamente recomendável, porém lembre-se que apenas assistir vídeos passivamente
não é suficiente para aprender Matemática, da mesma forma que, assistir atletas olímpicos não nos torna
atletas também.

Alguns recursos \textit{online} que podem ser uteis:
\begin{itemize}[topsep=-0.7\baselineskip]
    \item \href{http://www.wolframalpha.com}{www.wolframalpha.com} \\
         WolframAlpha é um mecanismo de conhecimento computacional
         que é capaz de fazer muitos cálculos.

    \item \href{http://pt.khanacademy.org}{pt.khanacademy.org} \\
        A Khan Academy é uma ONG educacional criada e sustentada por Sal Khan.
        Com a missão de fornecer educação de alta qualidade para qualquer um, em qualquer lugar,
        oferece uma coleção grátis de vídeos de matemática, medicina e saúde, economia e finanças,
        física, química, biologia, ciência da computação, entre outras matérias.

    \item \href{http://www.mathway.com/pt}{www.mathway.com/pt} \\
        A Mathway fornece aos alunos ferramentas para compreender e
        resolver problemas matemáticos.
        As resoluções são apresentadas passo a passo.

\end{itemize}

Nas seções desse capítulo indicamos alguns vídeos e atividades relacionados com os conteúdos de cada
capítulo dessa apostila. Esses conteúdos podem ser muito úteis como auxilio no estudo de cada tópico.
Note, porém, que a organização ou ordem dos tópicos varia de curso para curso. Algumas vezes notações
também variam e em casos extremos definições distintas podem ser empregadas.

%------------------------------------------------------------------------------%
\section{Integração}

As principais referências para esse tópico são:
\begin{listrefs}
  \item Stewart,\, Cálculo Vol. 1 \cite{stewart_1} -- Capítulo 5, Seções 5-1 até 5-4
  \item Thomas,    Cálculo Vol. 1 \cite{thomas_1}  -- Capítulo 5, Seções 5-1 até 5-4; Capítulo 8, Seção 8.8
\end{listrefs}

Aulas gravadas pelo professor Jonathas Douglas Santos de Oliveira:
\begin{listlinks}
  \item \href{https://youtu.be/8tyJhxZfiSY}{Revisão de Primitivas e Integrais indefinidas}
  \item \href{https://youtu.be/hD18trWiGEs}{Integrais definidas: A ideia de Área sob uma curva}
  \item \href{https://youtu.be/SOW8gYBs98c}{O Teorema Fundamental do Cálculo}
  \item \href{https://youtu.be/tYQ5nwtuFUQ}{Integrais Impróprias}
\end{listlinks}

Aulas e exercícios do Khan Accademy:
\begin{listlinks}
  \item \href{https://pt.khanacademy.org/math/integral-calculus/ic-integration/ic-integral-calc-intro/v/introduction-to-integral-calculus}{Introdução ao cálculo integral}
  \item \href{https://pt.khanacademy.org/math/integral-calculus/ic-integration/ic-riemann-sums/v/simple-riemann-approximation-using-rectangles}{Introdução à aproximação de Riemann}
  \item \href{https://pt.khanacademy.org/math/integral-calculus/ic-integration/ic-summation-notation/v/sigma-notation-sum}{Revisão da notação de somatório}
  \item \href{https://pt.khanacademy.org/math/integral-calculus/ic-integration/ic-riemann-sums-summation-notation/v/generalizing-a-left-riemann-sum-with-equally-spaced-rectangles}{Somas de Riemann em notação de somatório}
  \item \href{https://pt.khanacademy.org/math/integral-calculus/ic-integration/ic-definite-integral-definition/v/riemann-sums-and-integrals}{Definindo integrais com somas de Riemann}
  \item \href{https://pt.khanacademy.org/math/integral-calculus/ic-integration/ic-ftc-part-1/v/fundamental-theorem-of-calculus}{Teorema fundamental do cálculo e funções de acumulação}
  \item \href{https://pt.khanacademy.org/math/integral-calculus/ic-integration/ic-connect-function-antiderivative/v/interpreting-behavior-of-antiderivative}{Interpretação do comportamento de funções de acumulação}
  \item \href{https://pt.khanacademy.org/math/integral-calculus/ic-integration/ic-integral-prop/v/negative-definite-integrals}{Propriedades de integrais definidas}
  \item \href{https://pt.khanacademy.org/math/integral-calculus/ic-integration/ic-ftc-part-2/v/connecting-the-first-and-second-fundamental-theorems-of-calculus}{Teorema fundamental do cálculo e integrais definidas}
  \item \href{https://pt.khanacademy.org/math/integral-calculus/ic-integration/ic-reverse-power-rule/v/indefinite-integrals-of-x-raised-to-a-power}{Regra da potência reversa}
  \item \href{https://pt.khanacademy.org/math/integral-calculus/ic-integration/ic-common-indefinite-integrals/v/antiderivative-of-x-1}{Integrais indefinidas de funções comuns}
  \item \href{https://pt.khanacademy.org/math/integral-calculus/ic-integration/ic-common-definite-integrals/v/reverse-power-rule-for-definite-integrals}{Integrais definidas de funções comuns}
  \item \href{https://pt.khanacademy.org/math/integral-calculus/ic-integration/ic-improper-integrals/v/introduction-to-improper-integrals}{Integrais impróprias}
\end{listlinks}

%------------------------------------------------------------------------------%
\section{Áreas e Volumes}

As principais referências para esse tópico são:
\begin{listrefs}
  \item Stewart,\, Cálculo Vol. 1 \cite{stewart_1} -- Capítulo 6, Seções 6-1 até 6-3
  \item Thomas,    Cálculo Vol. 1 \cite{thomas_1}  -- Capítulo 5, Seção 5-6; Capítulo 6, Secções 6-1 e 6-2
\end{listrefs}

Aulas gravadas pelo professor Jonathas Douglas Santos de Oliveira:
\begin{listlinks}
  \item \href{https://youtu.be/LL2d3vOB31s}{Área Entre Curvas}
\end{listlinks}

%------------------------------------------------------------------------------%
\section{Técnicas de Integração}

As principais referências para esse tópico são:
\begin{listrefs}
  \item Stewart,\, Cálculo Vol. 1 \cite{stewart_1} -- Capítulo 5, Seção 5-5; Capítulo 7, Seções 7-1 até 7-5
  \item Thomas,    Cálculo Vol. 1 \cite{thomas_1}  -- Capítulo 8, Seções 8-1 até 8-5
\end{listrefs}

Aulas gravadas pelo professor Jonathas Douglas Santos de Oliveira:
\begin{listlinks}
  \item \href{https://youtu.be/aqnCC7TkHhI}{Substituição Simples}
  \item \href{https://youtu.be/a\_qvMOmJu0Q}{Integração por partes}
  \item \href{https://youtu.be/ijl8718-Xog}{Integrais Trigonométricas}
  \item \href{https://youtu.be/V-t\_HRuvw8Q}{Substituição Trigonométrica}
  \item \href{https://youtu.be/ztgPS7RVG7Y}{Integração por frações parciais}
  \item \href{https://youtu.be/iJG\_rojVgVs}{Exercícios: Técnicas de Integração}
\end{listlinks}

Aulas e exercícios do Khan Accademy:
\begin{listlinks}
  \item \href{https://pt.khanacademy.org/math/integral-calculus/ic-integration/ic-u-sub/v/u-substitution}{Integração por substituição}
  \item \href{https://pt.khanacademy.org/math/integral-calculus/ic-integration/ic-long-div/v/integral-partial-fraction}{Integração usando divisão longa e o método de completar quadrados}
  \item \href{https://pt.khanacademy.org/math/integral-calculus/ic-integration/ic-integration-with-trig-identities/v/using-trig-identity-to-use-u-substitution}{Integração usando identidades trigonométricas}
  \item \href{https://pt.khanacademy.org/math/integral-calculus/ic-integration/ic-trig-substitution/v/introduction-to-trigonometric-substitution}{Substituição trigonométrica}
  \item \href{https://pt.khanacademy.org/math/integral-calculus/ic-integration/ic-integration-by-parts/v/deriving-integration-by-parts-formula}{Integração por partes}
  \item \href{https://pt.khanacademy.org/math/integral-calculus/ic-integration/ic-partial-frac/v/integration-with-partial-fractions}{Integração usando frações parciais lineares}
\end{listlinks}

%------------------------------------------------------------------------------%
\section{Sequências Numéricas}

As principais referências para esse tópico são:
\begin{listrefs}
  \item Stewart,\, Cálculo Vol. 2 \cite{stewart_2} -- Capítulo 11, Seção 11-1
  \item Thomas,    Cálculo Vol. 2 \cite{thomas_2}  -- Capítulo 10, Seção 10-1
\end{listrefs}

Aulas gravadas pelo professor J. L. Acebal:
\begin{listlinks}
  \item \href{https://youtu.be/7A1s_n1Tpmk}{Aula 1 - Motivação para Sequências}
  \item \href{https://youtu.be/9OPgOfeZW0o}{Aula 2 - Sequências e Definições}
  \item \href{https://youtu.be/MqU6PqxAOHc}{Aula 3 - Sequências: limites e propriedades (Parte 1)}
  \item \href{https://youtu.be/GddAFc4l7CI}{Aula 3 - Sequências: limites e propriedades (Parte 2)}
\end{listlinks}

Aulas e exercícios do Khan Accademy:
\begin{listlinks}
  \item \href{https://pt.khanacademy.org/math/calculus-home/series-calc/sequences-calc/v/explicit-and-recursive-definitions-of-sequences}{Revisão sobre progressões}
  \item \href{https://pt.khanacademy.org/math/calculus-home/series-calc/seq-conv-diverg-calc/v/convergent-and-divergent-sequences       }{Sequências infinitas     }
\end{listlinks}

Aulas da disciplina Cálculo IV para a Engenharia (MAT-2456)
ministrada no segundo semestre de 2014 pelo Prof. Claudio Possani da USP:
\href{https://www.youtube.com/playlist?list=PLxI8Can9yAHeOiMYCBlkyCALloROQ58OY}%
     {MAT2456 - Cálculo Diferencial e Integral para Engenharia IV}.
\begin{listlinks}
  \item \href{https://www.youtube.com/watch?v=BrufXJ2AhNo&list=PLxI8Can9yAHeOiMYCBlkyCALloROQ58OY&index=2&t=0s}{Aula 1 - Sequências Numéricas I  - Parte 1 de 4}
  \item \href{https://www.youtube.com/watch?v=-WXsYcWLr30&list=PLxI8Can9yAHeOiMYCBlkyCALloROQ58OY&index=3&t=0s}{Aula 1 - Sequências Numéricas I  - Parte 2 de 4}
  \item \href{https://www.youtube.com/watch?v=j8bd446q2wk&list=PLxI8Can9yAHeOiMYCBlkyCALloROQ58OY&index=4&t=0s}{Aula 1 - Sequências Numéricas I  - Parte 3 de 4}
  \item \href{https://www.youtube.com/watch?v=767sM0V6Ss4&list=PLxI8Can9yAHeOiMYCBlkyCALloROQ58OY&index=5&t=0s}{Aula 1 - Sequências Numéricas I  - Parte 4 de 4}
  \item \href{https://www.youtube.com/watch?v=lUfv3bspiGw&list=PLxI8Can9yAHeOiMYCBlkyCALloROQ58OY&index=6&t=0s}{Aula 2 - Sequências Numéricas II - Parte 1 de 3}
  \item \href{https://www.youtube.com/watch?v=LxjSbzWipak&list=PLxI8Can9yAHeOiMYCBlkyCALloROQ58OY&index=7&t=0s}{Aula 2 - Sequências Numéricas II - Parte 2 de 3}
  \item \href{https://www.youtube.com/watch?v=hSuUyOXvfvk&list=PLxI8Can9yAHeOiMYCBlkyCALloROQ58OY&index=8&t=0s}{Aula 2 - Sequências Numéricas II - Parte 3 de 3}
\end{listlinks}

%------------------------------------------------------------------------------%
\section{Séries Numéricas}

As principais referências para esse tópico são:
\begin{listrefs}
  \item Stewart,\, Cálculo Vol. 2 \cite{stewart_2} -- Capítulo 11, Seções 11-2 até 11-7
  \item Thomas,    Cálculo Vol. 2 \cite{thomas_2}  -- Capítulo 10, Seções 10-2 até 10-6
\end{listrefs}

Aulas gravadas pelo professor J. L. Acebal:
\begin{listlinks}
  \item \href{https://youtu.be/-JzlrLOZrGc}{Aula 4 - Sequências de Somas Parciais e Séries}
  \item \href{https://youtu.be/DAer_v9z9cs}{Aula 5 - Séries importantes, Teste de Divergência}
  \item \href{https://youtu.be/Esl6Ot0Ewe8}{Aula 6 - Testes de Comparação}
  \item \href{https://youtu.be/bqM1tT5U3n4}{Aula 7 - Séries - Testes da Integral, da Razão e da Raiz}
  \item \href{https://youtu.be/4KRHu_4eiMc}{Aula 8 - Séries de Termos Positivos e Negativos - Convergência Absoluta e Convergência Condicional}
\end{listlinks}

Aulas e exercícios do Khan Accademy:
\begin{listlinks}
  \item \href{https://pt.khanacademy.org/math/calculus-home/series-calc/series-calculus/v/sigma-notation-sum                            }{Revisão de séries                   }
  \item \href{https://pt.khanacademy.org/math/calculus-home/series-calc/finite-geometric-series-calc/v/series-as-sum-of-sequence        }{Séries geométricas finitas          }
  \item \href{https://pt.khanacademy.org/math/calculus-home/series-calc/partial-sums-calc/v/partial-sum-notation                        }{Somas parciais                      }
  \item \href{https://pt.khanacademy.org/math/calculus-home/series-calc/geo-series-calc/v/evaluating-infinite-geometric-series          }{Séries geométricas infinitas        }
  \item \href{https://pt.khanacademy.org/math/calculus-home/series-calc/series-basics-challenge/v/telescoping-series                    }{Desafio de noções básicas de série  }
  \item \href{https://pt.khanacademy.org/math/calculus-home/series-calc/basic-convergence-tests-calc/v/divergence-test                  }{Testes de convergência básica       }
  \item \href{https://pt.khanacademy.org/math/calculus-home/series-calc/convergence-divergence-tests-calc/v/comparison-test             }{Testes de comparação                }
  \item \href{https://pt.khanacademy.org/math/calculus-home/series-calc/ratio-and-alternating-series-tests-calc/v/ratio-test-convergence}{Testes da razão e da série alternada}
  \item \href{https://pt.khanacademy.org/math/calculus-home/series-calc/estimating-infinite-series-calc/v/integrals-estimating-p-series }{Como estimar séries infinitas       }
\end{listlinks}

Aulas da disciplina Cálculo IV para a Engenharia (MAT-2456)
ministrada no segundo semestre de 2014 pelo Prof. Claudio Possani da USP:
\href{https://www.youtube.com/playlist?list=PLxI8Can9yAHeOiMYCBlkyCALloROQ58OY}%
     {MAT2456 - Cálculo Diferencial e Integral para Engenharia IV}.

Séries numéricas:
\begin{listlinks}
  \item \href{https://www.youtube.com/watch?v=8MjfOpAINUk&list=PLxI8Can9yAHeOiMYCBlkyCALloROQ58OY&index=9&t=0s}{Aula 2 - Séries Numéricas - Conceitos Básicos}
\end{listlinks}

Critérios de convergência:
\begin{listlinks}
  \item \href{https://www.youtube.com/watch?v=d_U9WfFXTXY&list=PLxI8Can9yAHeOiMYCBlkyCALloROQ58OY&index=10&t=0s}{Aula 3 - Critérios de Convergência - Parte 1 de 8}
  \item \href{https://www.youtube.com/watch?v=IqnYmacgmpY&list=PLxI8Can9yAHeOiMYCBlkyCALloROQ58OY&index=11&t=0s}{Aula 3 - Critérios de Convergência - Parte 2 de 8}
  \item \href{https://www.youtube.com/watch?v=jC9MWU3VrsU&list=PLxI8Can9yAHeOiMYCBlkyCALloROQ58OY&index=12&t=0s}{Aula 3 - Critérios de Convergência - Parte 3 de 8}
  \item \href{https://www.youtube.com/watch?v=UOoV1FX5oO8&list=PLxI8Can9yAHeOiMYCBlkyCALloROQ58OY&index=13&t=0s}{Aula 3 - Critérios de Convergência - Parte 4 de 8}
  \item \href{https://www.youtube.com/watch?v=F_Mx__unOeM&list=PLxI8Can9yAHeOiMYCBlkyCALloROQ58OY&index=14&t=0s}{Aula 4 - Critérios de Convergência - Parte 5 de 8}
  \item \href{https://www.youtube.com/watch?v=XZeLekfACKs&list=PLxI8Can9yAHeOiMYCBlkyCALloROQ58OY&index=15&t=0s}{Aula 4 - Critérios de Convergência - Parte 6 de 8}
  \item \href{https://www.youtube.com/watch?v=o0fS-j0BYqg&list=PLxI8Can9yAHeOiMYCBlkyCALloROQ58OY&index=16&t=0s}{Aula 4 - Critérios de Convergência - Parte 7 de 8}
  \item \href{https://www.youtube.com/watch?v=5x24l90oZFk&list=PLxI8Can9yAHeOiMYCBlkyCALloROQ58OY&index=17&t=0s}{Aula 4 - Critérios de Convergência - Parte 8 de 8}
\end{listlinks}

Convergência absoluta e condicional:
\begin{listlinks}
  \item \href{https://www.youtube.com/watch?v=PPu7D3bIwsc&list=PLxI8Can9yAHeOiMYCBlkyCALloROQ58OY&index=18&t=0s}{Aula 5 - Convergência Condicional ou Absoluta I}
  \item \href{https://www.youtube.com/watch?v=UGD6c-3whJA&list=PLxI8Can9yAHeOiMYCBlkyCALloROQ58OY&index=19&t=0s}{Aula 5 - Convergência Condicional ou Absoluta II}
\end{listlinks}

Aulas do curso de Cálculo III (MA311) na Unicamp ministrado pela Professora Ketty A. de Rezende:
\href{https://www.youtube.com/playlist?list=PLFBA21F349930F92F}{Cursos Unicamp - Cálculo III}

Sequências e séries Numéricas:
\begin{listlinks}
  \item \href{https://www.youtube.com/watch?v=JiWKLACcRME&list=PLFBA21F349930F92F&index=38&t=0s}{Séries Numéricas; Testes de Convergência - Parte 1}
  \item \href{https://www.youtube.com/watch?v=AVUjsHZ9Um4&list=PLFBA21F349930F92F&index=39&t=0s}{Séries Numéricas; Testes de Convergência - Parte 2}
  \item \href{https://www.youtube.com/watch?v=W9BCFKULvaU&list=PLFBA21F349930F92F&index=40&t=0s}{Testes de Convergência e das Séries Alternadas - Parte 1}
  \item \href{https://www.youtube.com/watch?v=cb0apXI5SaM&list=PLFBA21F349930F92F&index=41&t=0s}{Testes de Convergência e das Séries Alternadas - Parte 2}
\end{listlinks}

%------------------------------------------------------------------------------%
\section{Séries de Potências e de Taylor}

As principais referências para esse tópico são:
\begin{listrefs}
  \item Stewart, Cálculo Vol. 2 \cite{stewart_2} -- Capítulo 11, Seções 11-8 até 11-11
  \item Thomas,  Cálculo Vol. 2 \cite{thomas_2}  -- Capítulo 10, Seções 10-7 até 10-10
\end{listrefs}

Aulas gravadas pelo professor J. L. Acebal:
\begin{listlinks}
  \item \href{https://youtu.be/ph9g86DY2Ng}{Aula  9 - Séries de Potências - Raio de Convergência}
  \item \href{https://youtu.be/uSZmRjUjXMc}{Aula 10 - Séries de Potências - Propriedades Diferenciais e Algébricas}
  \item \href{https://youtu.be/W68BrWELpLw}{Aula 11 - Séries de Potências - Coeficientes de Taylor}
  \item \href{https://youtu.be/i5gxxsaRSdk}{Aula 12 - Séries de Potências - Polinômio de Taylor, resto de Taylor, erros de aproximação, linearização e convergência}
  \item \href{https://youtu.be/LCR2unh7zLU}{Aula 13 - Resolução de EDOs por Séries Potências}
\end{listlinks}

Aulas e exercícios do Khan Accademy:
\begin{listlinks}
  \item \href{https://pt.khanacademy.org/math/calculus-home/series-calc/power-series-intro-calc/v/power-series-radius-interval-convergence}{Introdução à série de potências}
  \item \href{https://pt.khanacademy.org/math/calculus-home/series-calc/taylor-series-calc/v/maclaurin-and-taylor-series-intuition        }{Introdução aos polinômios de Taylor e Maclaurin}
  \item \href{https://pt.khanacademy.org/math/calculus-home/series-calc/maclaurin-taylor-calc/v/cosine-taylor-series-at-0-maclaurin       }{Série de Maclaurin de $\sin(x)$, $\cos(x)$, e $e^x$}
  \item \href{https://pt.khanacademy.org/math/calculus-home/series-calc/power-series-calc/v/function-as-a-geometric-series                }{Representação de função de série de potência}
  \item \href{https://pt.khanacademy.org/math/calculus-home/series-calc/challenge-exercises-series-calc/e/convergence-tests-challenge     }{Desafios de séries}
\end{listlinks}

Aulas da disciplina Cálculo IV para a Engenharia (MAT-2456)
ministrada no segundo semestre de 2014 pelo Prof. Claudio Possani da USP:
\href{https://www.youtube.com/playlist?list=PLxI8Can9yAHeOiMYCBlkyCALloROQ58OY}{MAT2456 - Cálculo Diferencial e Integral para Engenharia IV}

Séries de Potências:
\begin{listlinks}
  \item \href{https://www.youtube.com/watch?v=M3459ZGJFd4&list=PLxI8Can9yAHeOiMYCBlkyCALloROQ58OY&index=20&t=0s}{Aula 5 - Séries de Potências - Introdução I}
  \item \href{https://www.youtube.com/watch?v=dEIytPdERKQ&list=PLxI8Can9yAHeOiMYCBlkyCALloROQ58OY&index=21&t=0s}{Aula 6 - Séries de Potências - Introdução II}
  \item \href{https://www.youtube.com/watch?v=EigL60vEwCc&list=PLxI8Can9yAHeOiMYCBlkyCALloROQ58OY&index=22&t=0s}{Aula 6 - Séries de Potências - Parte 1 de 5}
  \item \href{https://www.youtube.com/watch?v=MFO89oazdPQ&list=PLxI8Can9yAHeOiMYCBlkyCALloROQ58OY&index=23&t=0s}{Aula 6 - Séries de Potências - Parte 2 de 5 (Continua depois da Aula 7)}
  \item \href{https://www.youtube.com/watch?v=HqnYPbD3xvk&list=PLxI8Can9yAHeOiMYCBlkyCALloROQ58OY&index=24&t=0s}{Aula 7 - Revisão e Aprofundamento - Séries de Potências - Parte 1 de 7 }
  \item \href{https://www.youtube.com/watch?v=w0L8c4MkDjE&list=PLxI8Can9yAHeOiMYCBlkyCALloROQ58OY&index=25&t=0s}{Aula 7 - Revisão e Aprofundamento - Séries de Potências - Parte 2 de 7 }
  \item \href{https://www.youtube.com/watch?v=uIX-OaOat_8&list=PLxI8Can9yAHeOiMYCBlkyCALloROQ58OY&index=26&t=0s}{Aula 7 - Revisão e Aprofundamento - Séries de Potências - Parte 3 de 7 }
  \item \href{https://www.youtube.com/watch?v=HYSzJ3VKNaw&list=PLxI8Can9yAHeOiMYCBlkyCALloROQ58OY&index=27&t=0s}{Aula 7 - Revisão e Aprofundamento - Séries de Potências - Parte 4 de 7 }
  \item \href{https://www.youtube.com/watch?v=Y9Prh7iYVcg&list=PLxI8Can9yAHeOiMYCBlkyCALloROQ58OY&index=28&t=0s}{Aula 7 - Revisão e Aprofundamento - Séries de Potências - Parte 5 de 7 }
  \item \href{https://www.youtube.com/watch?v=NCxoy1Ncw8A&list=PLxI8Can9yAHeOiMYCBlkyCALloROQ58OY&index=29&t=0s}{Aula 7 - Revisão e Aprofundamento - Séries de Potências - Parte 6 de 7 }
  \item \href{https://www.youtube.com/watch?v=FcXHCFG2hDw&list=PLxI8Can9yAHeOiMYCBlkyCALloROQ58OY&index=30&t=0s}{Aula 7 - Revisão e Aprofundamento - Séries de Potências - Parte 7 de 7 }
  \item \href{https://www.youtube.com/watch?v=ZkP8S6-2Ljo&list=PLxI8Can9yAHeOiMYCBlkyCALloROQ58OY&index=31&t=0s}{Aula 8 - Séries de Potências - Parte 3 de 5 (Continuação da Aula 6)}
  \item \href{https://www.youtube.com/watch?v=8ad2ly1clvc&list=PLxI8Can9yAHeOiMYCBlkyCALloROQ58OY&index=32&t=0s}{Aula 8 - Séries de Potências - Parte 4 de 5}
  \item \href{https://www.youtube.com/watch?v=Dmiud-Xoy0A&list=PLxI8Can9yAHeOiMYCBlkyCALloROQ58OY&index=33&t=0s}{Aula 8 - Séries de Potências - Parte 5 de 5}
\end{listlinks}

Raio de convergência,
derivação e integração termo-a-termo,
série de Taylor:
\begin{listlinks}
  \item \href{https://www.youtube.com/watch?v=XNg3c8Pm53U&list=PLxI8Can9yAHeOiMYCBlkyCALloROQ58OY&index=34&t=0s}{Aula 9  - Séries de Taylor - Parte 1 de 7}
  \item \href{https://www.youtube.com/watch?v=g8giCcjyl-8&list=PLxI8Can9yAHeOiMYCBlkyCALloROQ58OY&index=35&t=0s}{Aula 9  - Séries de Taylor - Parte 2 de 7}
  \item \href{https://www.youtube.com/watch?v=4Vd4W0qc68I&list=PLxI8Can9yAHeOiMYCBlkyCALloROQ58OY&index=36&t=0s}{Aula 9  - Séries de Taylor - Parte 3 de 7}
  \item \href{https://www.youtube.com/watch?v=CavnRpUt9bA&list=PLxI8Can9yAHeOiMYCBlkyCALloROQ58OY&index=37&t=0s}{Aula 9  - Séries de Taylor - Parte 4 de 7}
  \item \href{https://www.youtube.com/watch?v=mD0D0c5_9iE&list=PLxI8Can9yAHeOiMYCBlkyCALloROQ58OY&index=38&t=0s}{Aula 10 - Séries de Taylor - Parte 5 de 7}
  \item \href{https://www.youtube.com/watch?v=E72lVIEWkPg&list=PLxI8Can9yAHeOiMYCBlkyCALloROQ58OY&index=39&t=0s}{Aula 10 - Séries de Taylor - Parte 6 de 7}
  \item \href{https://www.youtube.com/watch?v=uTkn4KJPSr0&list=PLxI8Can9yAHeOiMYCBlkyCALloROQ58OY&index=40&t=0s}{Aula 10 - Séries de Taylor - Parte 7 de 7}
\end{listlinks}

Revisão e Aprofundamento:
\begin{listlinks}
  \item \href{https://www.youtube.com/watch?v=MSqsEiOUwWU&list=PLxI8Can9yAHeOiMYCBlkyCALloROQ58OY&index=49&t=0s}{Aula 13 - Exercícios de Séries de Potências}
  \item \href{https://www.youtube.com/watch?v=M3Ve6qPsY9I&list=PLxI8Can9yAHeOiMYCBlkyCALloROQ58OY&index=52&t=0s}{Aula 14 - Revisão e Aprofundamento - Séries de Taylor - Parte 1 de 4}
  \item \href{https://www.youtube.com/watch?v=ldOL7jfU9aI&list=PLxI8Can9yAHeOiMYCBlkyCALloROQ58OY&index=53&t=0s}{Aula 14 - Revisão e Aprofundamento - Séries de Taylor - Parte 2 de 4}
  \item \href{https://www.youtube.com/watch?v=hQfM2tF1VQg&list=PLxI8Can9yAHeOiMYCBlkyCALloROQ58OY&index=50&t=0s}{Aula 14 - Revisão e Aprofundamento - Séries de Taylor - Parte 3 de 4}
  \item \href{https://www.youtube.com/watch?v=sDUW26Yhmf8&list=PLxI8Can9yAHeOiMYCBlkyCALloROQ58OY&index=51&t=0s}{Aula 14 - Revisão e Aprofundamento - Séries de Taylor - Parte 4 de 4}
\end{listlinks}

Aulas do curso de Cálculo III (MA311) na Unicamp ministrado pela Professora Ketty A. de Rezende:
\href{https://www.youtube.com/playlist?list=PLFBA21F349930F92F}{Cursos Unicamp - Cálculo III}

Séries de Potências e de Taylor:
\begin{listlinks}
  \item \href{https://www.youtube.com/watch?v=w1M69ifBFX4&list=PLFBA21F349930F92F&index=42&t=0s}{Séries de Potências - Parte 1}
  \item \href{https://www.youtube.com/watch?v=PXmcPSLZxmQ&list=PLFBA21F349930F92F&index=43&t=0s}{Séries de Potências - Parte 2}
  \item \href{https://www.youtube.com/watch?v=8STPRjedg44&list=PLFBA21F349930F92F&index=44&t=0s}{Série de Potências em Ponto Ordinário - Parte 1}
  \item \href{https://www.youtube.com/watch?v=C1FRI5CRQIo&list=PLFBA21F349930F92F&index=45&t=0s}{Série de Potências em Ponto Ordinário - Parte 2}
\end{listlinks}

%------------------------------------------------------------------------------%
